% !TeX spellcheck = he_IL
%% ! TEX TS-program = xelatexmk
% above: commands appropriate for TeXShop in mac OS
\documentclass[12pt, a4paper]{article}
\usepackage{etoolbox}
\usepackage{bbm}
%\usepackage{pifont}
\usepackage{tikz}
%\input{preamble/UnitPreamble}
%%%%%%%%%%%%%%%%%%%%%%%%%%
%%
%%
%%  Exam Template
%%
%%
%%%%%%%%%%%%%%%%%%%%%%%%%%
\newtoggle{testreview}
\toggletrue{testreview}  
\togglefalse{testreview} % uncomment to make it ready for exam unit
\newtoggle{testsolution}
\toggletrue{testsolution}  
\togglefalse{testsolution} % comment to show solutions
\newcounter{version}
\setcounter{version}{2}  % version of multiple-choice

	





\usepackage{graphicx}

\usepackage{braket}
\usepackage[a4paper, top=3cm,left=2.5cm, right=2.3cm, bottom=2.7cm]{geometry}
\usepackage{amsmath, amssymb, array, multicol, enumitem}
\usepackage{pdfpages}
\usepackage{comment, graphicx}  
\usepackage{xparse}
\DeclareGraphicsExtensions{.png,.jpg}
\usepackage{pgffor}
\usepackage[colorlinks=true]{hyperref} 
\AtBeginDocument{%
	\hypersetup{pdfborder={0 0 0},
		urlcolor=blue}%
}

\usepackage{wasysym}
\usepackage{fontspec}
\usepackage{xunicode}
\usepackage{polyglossia}
\setdefaultlanguage{hebrew}
\setotherlanguage{english}
%\newfontfamily\hebrewfont[Script=Hebrew, Ligatures=TeX, Scale=1.1, WordSpace =1.05]{David}
%\newfontfamily\hebrewfont[Script=Hebrew, Ligatures=TeX, Scale=1.1, WordSpace=1.05]{David  CLM}
\newfontfamily\hebrewfont[Script=Hebrew]{Times New Roman}
%\newfontfamily\hebrewfont[Script=Hebrew]{Ezra SIL}

% Hebrew monospace font (must support Hebrew)
\newfontfamily\hebrewfonttt[Script=Hebrew]{FreeMono}

%http://tex.stackexchange.com/questions/141231/equation-number-enclosing-parentheses-misplaced-with-polyglossiahebrewamsmath
\makeatletter
%\def\maketag@@@#1{\hbox{\m@th\normalfont\LRE{#1}}}
\def\tagform@#1{\maketag@@@{(\ignorespaces#1\unskip)}}
\makeatother

%
% Hebrew fonts are two small 
%
\fontsize{14}{15}

%
% Package fro proper positioning of logo
%
\usepackage[absolute]{textpos}


% http://tex.stackexchange.com/questions/10589/fix-nested-section-numbers-in-rtl-languages-with-polyglossia/
\makeatletter
\renewcommand{\SepMark}[1]{\def\@SepMark{#1}}
\makeatother
\SepMark{.}
\robustify\hebrewnumeral
\robustify\Hebrewnumeral
\robustify\Hebrewnumeralfinal


% from polyglossia docs
%\gappto\captionshebrew{%
\def\urladdrname{{\itshape כתובת רשת}}
\def\emailaddrname{{\itshape כתובת דואל}}
%}

% https://tex.stackexchange.com/questions/176661/polyglossia-arabic-text-english-bibliography
\makeatletter
\let\orig@MakeUppercase \MakeUppercase
\makeatother
\providecommand{\MakeUppercase}[1]{#1}




%%%%%%%%%%%%%%%%%%%%%%%%%%
%%
%%
%%   User Commands
%%
%%
%%%%%%%%%%%%%%%%%%%%%%%%%%
\newcommand{\heb}[1]{\begin{small}\R{#1}\end{small}}
\newcommand{\hebt}[1]{\begin{scriptsize}\R{#1}\end{scriptsize}}
\newcommand{\cc}{\mathbb{C}}
\newcommand{\nn}{\mathbb{N}}
\newcommand{\zz}{\mathbb{Z}}
%\renewcommand{\phi}{\varphi}
\newcommand{\abs}[1]{\left|#1\right|}
\newcommand{\lr}[1]{\left(#1\right)}


\newcommand{\ff}{\mathbb{F}}
\newcommand{\FF}{\mathbb{F}}
\newcommand{\mtn}{m\times n}
\newcommand{\longr}{\longrightarrow}
\newcommand{\longrl}{\longleftrightarrow}
\newcommand{\ol}{\overline}
\newcommand{\ntn}{n\times n}
\newcommand{\rr}{\mathbb{R}}
\newcommand{\QQ}{\mathbb{Q}}
\newcommand{\RR}{\mathbb{R}}
\newcommand{\im}{\mathrm{Im}}

\newcommand{\CC}{\mathbb{C}}
\newcommand{\ZZ}{\mathbb{Z}}
\newcommand{\qq}{\mathbb{Q}}
\DeclareMathOperator{\vol}{vol}

\DeclareMathOperator{\Adj}{Adj}
\DeclareMathOperator{\deter}{det}
\DeclareMathOperator{\Img}{Im}
\DeclareMathOperator{\Ker}{Ker}
\DeclareMathOperator{\rank}{rank}
\DeclareMathOperator{\Span}{Span}
\DeclareMathOperator{\tr}{tr}


\newcommand{\cB}{\mathcal{B}}
\newcommand{\cC}{\mathcal{C}}
\newcommand{\cR}{\mathcal{R}}
\newcommand{\cD}{\mathcal{D}}
\newcommand{\calE}{\mathcal{E}}


%%%% new for complex functions

\newcommand{\Arg}{\mathrm{Arg}}
\newcommand{\wt}{\widetilde}
\newcommand{\wh}{\widehat}
\newcommand{\re}{\mathrm{Re}}
\newcommand{\goes}{\longrightarrow}
\newcommand{\cO}{\mathcal{O}}
\newcommand{\Log}{\mathrm{Log}}


\newcommand{\twopartdef}[4]
{
	\left\{
		\begin{array}{ll}
			#1 & \mbox{   } #2 \\
			#3 & \mbox{   } #4
		\end{array}
	\right.
}

\newcommand{\solutionword}{\textbf{פתרון:}}
\newcommand{\points}[1]{(#1 נק') }
\iftoggle{testsolution}{%
	\newcommand{\sols}[1]{\solutionword$\;\;$ #1$\square$}
	\newcommand{\spacing}{0cm}
	\newcommand{\questionspacing}{2.5cm}
	\newcommand{\extrapages}{}
}{ %
%\newcommand{\spacing}{12cm}
%\newcommand{\halfspacing}{0.5cm}	
\newcommand{\spacing}{24cm}
\newcommand{\questionspacing}{\spacing}
\newcommand{\sols}[1]{}
\newcommand{\extrapages}{%
	\newpage
	\noindent
	דף זה ריק במכוון ומיועד לשימוש הסטודנט/ית.
	\newpage
	\noindent
	דף זה ריק במכוון ומיועד לשימוש הסטודנט/ית.
}
\iftoggle{testreview}{%
\renewcommand{\spacing}{1cm}
%\newcommand{\questionspacing}{\spacing}
\renewcommand{\extrapages}{}	
}{}	
}


\DeclareMathOperator{\row}{Row}

\DeclareMathOperator{\col}{Col}


%%%%%%%%%%%%%%%%%%%%%%%%
%%
%%   BEGIN DOCUMENT
%%
%%%%%%%%%%%%%%%%%%%%%%%%

\begin{document}

\selectlanguage{hebrew}


\newcommand{\semester}{%
	סמסטר א' תשפ״ו
}
\newcommand{\course}{%
 אלגברה לינארית להנדסה 
}
\newcommand{\coursenumber}{9321}	
\newcommand{\lecturers}{%
	מרצים: ד״ר י.מייזל, ד״ר נ.גולקו, ד״ר נ.מאיר, ד״ר נ.איידלשטין, ד״ר מ.פורת, ד״ר מ.זיו-אב, ד״ר מ.ברשדסקי
}

\newcommand{\testdate}{{3.2.26}}
\iftoggle{testsolution}{%
	\newcommand{\testtype}{%
		פתרונות למבחן מועד א'
	}
}{%
	\newcommand{\testtype}{%
	\boxed{\text{
	מועד א'
    -
    פתרון מלא
	}	}
}}

\newcommand{\duration}{%
	3 שעות.
}
\newcommand{\helpmaterial}{%
מחשבון פשוט בלבד.
}



\thispagestyle{empty}

\begin{textblock*}{9cm}[0.5,0.5](90mm,3cm)
	\includegraphics[width=5cm]{logo.jpg}
\end{textblock*}	

\noindent

\begin{center}
אוניברסיטת בן-גוריון בנגב - היחידה להוראת קורסי יסוד - %
\smallskip
\semester \\
{\large \bf \course - 214.1.\coursenumber}\\
\smallskip
\lecturers\\
\smallskip
\smallskip
\smallskip
{\Large \testtype }\\
\end{center}

\noindent
\textbf{%
תאריך:
} \testdate \quad 
\textbf{משך הבחינה: 
} \duration \quad
\textbf{
חומר עזר: 
}
\helpmaterial \\
דף הנוסחאות מצורף למבחן. אין להכניס כל חומר עזר נוסף.
\subsection*{חלק א': שאלות רב ברירה}
\begin{itemize}
    \item בחלק א' של המבחן יש 8 שאלות רב-ברירה, ומשקל כל שאלה 5 נקודות.
    \item לכל שאלה יש תשובה אחת נכונה בדיוק. יש לסמן תשובה אחת לכל שאלה \textbf{בטבלה בדף זה}, אין צורך בנימוק.
    \item שימו לב להנחיות מילוי הטבלה בעמוד הבא.
\end{itemize}

\begin{center}
\begin{tabular}{|@{}c@{}||@{}c@{}||@{}c@{}||@{}c@{}||@{}c@{}||@{}c@{}||@{}c@{}||@{}c@{}|}
\hline
1 & 2 &  3 & 4 & 5 & 6 & 7 & 8 \\
\hline
\begin{tabular}{cccc}
\\\textcircled{א}\\\\
\textcircled{ב}\\\\
\textcircled{ג}\\\\
\textcircled{ד}\\\\
\end{tabular}
&
\begin{tabular}{cccc}
\\\textcircled{א}\\\\
\textcircled{ב}\\\\
\textcircled{ג}\\\\
\textcircled{ד}\\\\
\end{tabular}

&
\begin{tabular}{cccc}
\\\textcircled{א}\\\\
\textcircled{ב}\\\\
\textcircled{ג}\\\\
\textcircled{ד}\\\\
\end{tabular}
&
\begin{tabular}{cccc}
\\\textcircled{א}\\\\
\textcircled{ב}\\\\
\textcircled{ג}\\\\
\textcircled{ד}\\\\
\end{tabular}
&
\begin{tabular}{cccc}
\\\textcircled{א}\\\\
\textcircled{ב}\\\\
\textcircled{ג}\\\\
\textcircled{ד}\\\\
\end{tabular}
&
\begin{tabular}{cccc}
\\\textcircled{א}\\\\
\textcircled{ב}\\\\
\textcircled{ג}\\\\
\textcircled{ד}\\\\
\end{tabular}
&
\begin{tabular}{cccc}
\\\textcircled{א}\\\\
\textcircled{ב}\\\\
\textcircled{ג}\\\\
\textcircled{ד}\\\\
\end{tabular}
&
\begin{tabular}{cccc}
\\\textcircled{א}\\\\
\textcircled{ב}\\\\
\textcircled{ג}\\\\
\textcircled{ד}\\\\
\end{tabular}
\\
\hline 
\begin{tikzpicture}
  \draw[dotted] (0,0) circle (0.8cm);
\end{tikzpicture}
& 
 \begin{tikzpicture}
  \draw[dotted] (0,0) circle (0.8cm);
\end{tikzpicture}
 &
 \begin{tikzpicture}
  \draw[dotted] (0,0) circle (0.8cm);
\end{tikzpicture}
 &
 \begin{tikzpicture}
  \draw[dotted] (0,0) circle (0.8cm);
\end{tikzpicture}
 &
 \begin{tikzpicture}
  \draw[dotted] (0,0) circle (0.8cm);
\end{tikzpicture}
 &
 \begin{tikzpicture}
  \draw[dotted] (0,0) circle (0.8cm);
\end{tikzpicture}
 &
 \begin{tikzpicture}
  \draw[dotted] (0,0) circle (0.8cm);
\end{tikzpicture}
 &
 \begin{tikzpicture}
  \draw[dotted] (0,0) circle (0.8cm);
\end{tikzpicture}
   \\
 \hline
\end{tabular}
\end{center}


\subsection*{חלק ב': שאלות פתוחות}
\begin{itemize}
    \item יש לבחור בדיוק 3 מתוך 4 שאלות פתוחות. משקל כל שאלה 20 נקודות.

    \item יש לסמן את השאלות שנבחרו בטבלה בדף זה. שאלות שלא סומנו לא תבדקנה!
    \item יש לנמק את התשובות בחלק ב' ולכתוב את כל החישובים הרלוונטים.
    \item יש לענות על כל השאלות בחלק זה על טופס המבחן במקום המיועד אחרי כל שאלה.   
    
    \underline{
    ניתן גם לכתוב מצדו השני של הדף. 
    }
    מחברת הטיוטה לא תבדק!
\end{itemize}
\bigskip

 בחלק ב' יש לבדוק את השאלות הבאות. סמנו את התשובות המתאימות --  3 שאלות לכל היותר!


\begin{center}
\setlength{\arrayrulewidth}{0.6pt}
\renewcommand{\arraystretch}{1.4}
\begin{tabular}{|c|c|c|c|}
\hline
\textcircled{1} & \textcircled{2} & \textcircled{3} & \textcircled{4}\\
\hline
\end{tabular}
\end{center}

\begin{center}
\huge{\textbf{בהצלחה!}}
\end{center}
\pagebreak
\subsection*{הוראות למילוי הטבלה}
\begin{itemize}
    \item[$\bullet$] יש להשתמש בעט שחור או כחול. 
    \item[$\bullet$] יש לסמן תשובה אחת בלבד ע"י השחרת העיגול המתאים.
    \item[$\bullet$]
    בשאלה בה לא יושחר אף עיגול,
    הניקוד יהיה 0.
    \item[$\bullet$] בשאלה בה יושחרו יותר מעיגול אחד, הניקוד יהיה 0.
.
    \item[$\bullet$] כדי לתקן סימון, יש לסמן X
    על התשובה הלא נכונה ולהשחיר את התשובה הנכונה.

    \item[$\bullet$]
    \textbf{שימו לב:}
    לא ניתן לבחור סעיף שסומן עליו X
    לכן הקפידו לסמן את תשובתכם כשאתם בטוחים בה.
\end{itemize}

\subsection*{דוגמאות}


\begin{center}
\begin{tabular}{|c||c||c||c||c||c||c||c||c||c|}
\hline
1 & 2 &  3 & 4 & 5 & 6 & 7 & 8 & 9 & 10 \\
\hline
\begin{tabular}{c}
\\\Large{\CIRCLE}\\\\
\ \  \hspace{-8pt}\raisebox{0pt}[0pt][0pt]{\Large\CIRCLE}%
\  \hspace{4pt}\llap{\raisebox{-4pt}[0pt][0pt]{ $\scalebox{3}{$\times$}$}}\hspace{-8pt}
\\\\
\textcircled{ג}\\\\
\textcircled{ד}\\\\
\end{tabular}
&
\begin{tabular}{cccc}
\\\textcircled{א}\\\\
\Large{\CIRCLE}\\\\
\textcircled{ג}\\\\
\textcircled{ד}\\\\
\end{tabular}
&
\begin{tabular}{cccc}
\\
\textcircled{א}
\\\\
\ \  \hspace{-8pt}\raisebox{0pt}[0pt][0pt]{\Large\CIRCLE}%
 \ \hspace{4pt}\llap{\raisebox{-4pt}[0pt][0pt]{ $\scalebox{3}{$\times$}$}}\hspace{-8pt}
\\\\
\Large{\CIRCLE}\\\\
\ \  \hspace{-8pt}\raisebox{0pt}[0pt][0pt]{\Large\CIRCLE}%
 \ \hspace{4pt}\llap{\raisebox{-4pt}[0pt][0pt]{ $\scalebox{3}{$\times$}$}}\hspace{-8pt}
\\\\
\end{tabular}
&
\begin{tabular}{cccc}
\\\textcircled{א}\\\\
\textcircled{ב}\\\\
\textcircled{ג}\\\\
\textcircled{ד}\\\\
\end{tabular}
&
\begin{tabular}{cccc}
\\\hspace{-8pt}\ \  \raisebox{0pt}[0pt][0pt]{\Large\CIRCLE}%
 \ \hspace{4pt}\llap{\raisebox{-4pt}[0pt][0pt]{ $\scalebox{3}{$\times$}$}}\hspace{-8pt}
\\\\
\textcircled{ב}\\\\
\textcircled{ג}\\\\
\textcircled{ד}\\\\
\end{tabular}
&
\begin{tabular}{cccc}
\\\hspace{-8pt}\ \  \raisebox{0pt}[0pt][0pt]{\Large\CIRCLE}%
 \ \hspace{4pt}\llap{\raisebox{-4pt}[0pt][0pt]{ $\scalebox{3}{$\times$}$}}\hspace{-8pt}
\\\\
\hspace{-8pt}\ \  \raisebox{0pt}[0pt][0pt]{\Large\CIRCLE}%
 \ \hspace{4pt}\llap{\raisebox{-4pt}[0pt][0pt]{ $\scalebox{3}{$\times$}$}}\hspace{-8pt}
\\\\
\textcircled{ג}\\\\
\textcircled{ד}\\\\
\end{tabular}
&
\begin{tabular}{cccc}
\\\Large{\CIRCLE}\\\\
\Large{\CIRCLE}\\\\
\textcircled{ג}\\\\
\textcircled{ד}\\\\
\end{tabular}
&
\begin{tabular}{cccc}
\\\Large{\CIRCLE}\\\\
\Large{\CIRCLE}\\\\
\Large{\CIRCLE}\\\\
\textcircled{ד}\\\\
\end{tabular}
&
\begin{tabular}{cccc}
\\\hspace{-8pt}\ \  \raisebox{0pt}[0pt][0pt]{\Large\CIRCLE}%
 \ \hspace{4pt}\llap{\raisebox{-4pt}[0pt][0pt]{ $\scalebox{3}{$\times$}$}}\hspace{-8pt}
\\\\
\Large{\CIRCLE}\\\\
\textcircled{ג}\\\\
\Large{\CIRCLE}\\\\
\end{tabular}
&
\begin{tabular}{cccc}
\\\hspace{-8pt}\ \  \raisebox{0pt}[0pt][0pt]{\Large\CIRCLE}%
 \ \hspace{4pt}\llap{\raisebox{-4pt}[0pt][0pt]{ $\scalebox{3}{$\times$}$}}\hspace{-8pt}
\\\\
\Large{\CIRCLE}
\\\\
\hspace{-8pt}\ \  \raisebox{0pt}[0pt][0pt]{\Large\CIRCLE}%
 \ \hspace{4pt}\llap{\raisebox{-4pt}[0pt][0pt]{ $\scalebox{3}{$\times$}$}}\hspace{-8pt}
\\\\
\Large{\CIRCLE}\\\\
\end{tabular}
\\
\hline
\end{tabular}
\ 

\ 

\begin{itemize}[label=$\bullet$]
    \item שאלות 1,2,3 עם סימונים תקינים עבור תשובות:
    \begin{enumerate}
        \item א
        \item ב
        \item ג
    \end{enumerate}
    \item בשאלות 4,5,6 חסר סימון והניקוד יהיה 0.
    \item בשאלות 7,8,9,10
    יש יותר מסימון אחד והניקוד יהיה 0
\end{itemize}


\end{center}


%%%
%%%
%%%   End of Header and Instructions
%%%
%%%

\pagebreak
%%%%%%%%%%%%%%%%%%%%%%%%%%%%%%%%%%%%%%
%%  Problems
%%%%%%%%%%%%%%%%%%%%%%%%%%%%%%%%%%%%%%


%%%%% שאלה 2 %%%%%%%

\section*{
חלק א' 
(40 נק')
להשחיר את התשובות בטבלה!}
\begin{enumerate}
\item\points{5}
יהיו
\[
v =
\begin{pmatrix}
1\\
1\\
1
\end{pmatrix}
\in \mathbb{C}^3
, \ 
B=\left\{
\begin{pmatrix}
1\\0\\0
\end{pmatrix},
\begin{pmatrix}
0\\ i\\0
\end{pmatrix},
\begin{pmatrix}
-1\\ a\\1
\end{pmatrix}
\right\}
\]
ונתון שוקטור הקואורדינטות של $v$ לפי הבסיס הסדור $B$ הוא
$
[v]_B = (2,1,1)
$.

מצאו את ערך הפרמטר \(a\).

\begin{enumerate}
\begin{multicols}{4}
\item $\boxed{1-i}$
\item \(-i\)
\item \(1\)
\item לא ניתן לדעת
\end{multicols}
\end{enumerate}

\noindent\rule{\textwidth-30pt}{0.4pt}

פתרון:
מהגדרת וקטור קואורדינטות נקבל כי
$$v=2\cdot \begin{pmatrix}
    1\\0\\0\\
\end{pmatrix}+1\cdot\begin{pmatrix}
    0\\i\\0\\
\end{pmatrix}+1\cdot\begin{pmatrix}
    -1\\a\\1\\
\end{pmatrix}=\begin{pmatrix}
    1\\i+a\\
    1\\
\end{pmatrix}$$
ולכן
$i+a=1$,
כלומר
$a=1-i$.

\noindent\rule{\textwidth-30pt}{0.4pt}

\item\points{5}
יהיו \(v_1,v_2\in\mathbb{R}^3\), כאשר
$
v_1=
(
1,2,2)$,
\ \ 
$\|v_2\|=1$.

מהם כל ערכי  \(\alpha\) עבורם הווקטורים
\((v_1+\alpha v_2),\quad (v_1-\alpha v_2)\) \ \ 
מאונכים זה לזה?

\begin{enumerate}
\begin{multicols}{4}
\item \(3\)
\item \(\pm \frac{3}{2}\)
\item \(\boxed{\pm 3}\)
\item \(\pm\sqrt{3}\)
\end{multicols}
\end{enumerate}

\noindent\rule{\textwidth-30pt}{0.4pt}

פתרון:
 בשביל לבדוק אם הם מאונכים, נחשב את המכפלה הסקלרית שלהם
 \begin{align*}
     (v_1+\alpha v_2)\bullet (v_1-\alpha v_2)=v_1\bullet v_1-\alpha v_1\bullet v_2+\alpha v_2\bullet v_1 -\alpha^2\cdot v_2\bullet v_2
 \end{align*}
 בגלל שמתקיים
 $v_1\bullet v_2=v_2\bullet v_1$
 ובגלל שמתקיים
 $v_1\bullet v_1=1\cdot 1+2\cdot 2+2\cdot 2=9$
 וגם
 $v_2\bullet v_2=\|v_2\|^2=1$,
 אז נסכם כי
 $ (v_1+\alpha v_2)\bullet (v_1-\alpha v_2)=9-\alpha^2$
 ולכן הוקטורים מאונכים אם ורק אם המכפלה הזו שווה 
 $0$,
 כלומר אם ורק אם
 $\alpha^2=9,$
 מה שאומר
 $\alpha=\pm3$.
 
\noindent\rule{\textwidth-30pt}{0.4pt}



%%%%% שאלה 2 %%%%%%%

\item\points{5}
נגדיר את $L$
להיות ישר החיתוך של שני המישורים הנתונים על ידי המשוואות הבאות:
\[
x-4z=1,\qquad
y-3z=-2
\]
איזה מישור מכיל את $L$?

\begin{enumerate}
\item $\Set{(x,y,z)\in \rr^3\ |\ x-y-z = 4}$
\item $\boxed{\Set{(x,y,z)\in \rr^3\ |\ 5x-3y-11z=11}}$
\item $\Set{(x,y,z)\in \rr^3\ |\ 2x+y=11}$
\item $\Set{(x,y,z)\in \rr^3\ |\ x-4y-z=9}$
\end{enumerate}

\noindent\rule{\textwidth-30pt}{0.4pt}


פתרון:
נרצה למצוא את החיתוך בין המישורים (שגילו לנו כי זה יוצא ישר): החיתוך נתון לנו על ידי
$$\{(x,y,z): x=1+4z, y=3z-2\}=\{(1+4z, 3z-2,z):z\in\rr\}$$
והרי זוהי הצגה פרמטרית של ישר העובר בנקודה
$(1,-2,0)$
עם וקטור כיוון
$(4,3,1)$.
הנה שתי דרכים לסיים את התרגיל:

$\bullet$
ניתן להציב את הנקודה
$(1,-2,0)$
ולראות כי שייכת רק לשניים מהמישורים. בנוסף, הכיוון
$(4,3,1)$
אמור להיות מאונך לנורמל של המישור, 
וזה אכן נכון כאשר הנורמל הוא
$(1,-1,-1)$,
או שהנורמל הוא
$(5,-3,-11)$,
אבל לא כאשר הוא 
$(1,-4,-1)$
או
$(2,1,0)$,
לכן שני מישורים נפסלים.
המישור היחיד שעומד בשני הקריטריונים הוא
המישור הנתון על ידי המשוואה
$5x-3y-11z=11$.

$\bullet$
 נציב את הנקודה הכללית מהישר שמצאנו בתוך כל אחת מהמשוואות של המישורים הנתונים, ונוודא כי אחת תמיד מתקיימת, כאשר האחרות לא.
 
\noindent\rule{\textwidth-30pt}{0.4pt}

\item\points{5}
תהי 
$A\in M_{4\times 4}(\rr)$
מטריצה הפיכה ותהי
\[
B=
\begin{pmatrix}
    1 & 1 & 1 & 1 \\
    1 & 2 & 2 & 2 \\
    1 & 2 & 3 & 3 \\
    1 & 2 & 3 & 4
\end{pmatrix}
\]
$\det(-2A^tBA^{-1})=$
\begin{enumerate}
\begin{multicols}{4}
    \item $-2$
    \item $4$
    \item $-8$
        \item $\boxed{16}$
\end{multicols}
\end{enumerate}


\noindent\rule{\textwidth-30pt}{0.4pt}


פתרון:
ראשית נשתמש בכללי הדטרמיננטה ונקבל
$$ֿ\det(-2A^tBA^{-1})=(-2)^4\cdot \det(A)\cdot \det(B)\det(A)^{-1}=16\cdot \det(B)$$
כעת נחשב את הדטרמיננטה של
$B$
וניעזר בביצוע פעולות שורה אלמנטריות
\begin{align*}
    \det(B)=\det\begin{pmatrix}
         1 & 1 & 1 & 1 \\
    1 & 2 & 2 & 2 \\
    1 & 2 & 3 & 3 \\
    1 & 2 & 3 & 4\\
    \end{pmatrix}
    &\underset{\begin{bmatrix}
        R_4-R_3\rightarrow R_4\\
    \end{bmatrix}}{=}
    \det\begin{pmatrix}
         1 & 1 & 1 & 1 \\
    1 & 2 & 2 & 2 \\
    1 & 2 & 3 & 3 \\
    0 & 0 & 0 & 1\\
    \end{pmatrix}
    \\&\underset{\begin{bmatrix}
        R_3-R_2\rightarrow R_3\\
        R_2-R_1\rightarrow R_2\\
    \end{bmatrix}}{=}
    \det\begin{pmatrix}
         1 & 1 & 1 & 1 \\
    0 & 1 & 1 & 1 \\
    0 & 0 & 1 & 1 \\
    0 & 0 & 0 & 1\\
    \end{pmatrix}=1
\end{align*}
ולסיכום נקבל כי התשובה הנכונה היא
$16$.

\noindent\rule{\textwidth-30pt}{0.4pt}

\pagebreak

\item\points{5}
תהי
\[A=\begin{pmatrix}
    a & 1 & 1 \\
    1 & a & 1 \\
    1 & 1 & a
\end{pmatrix}\]
עבור אילו ערכים של הפרמטר 
$a$
מתקיים 
$(4,-5,1)\in\row(A)$.
\begin{enumerate}
\begin{multicols}{2}
\item $a\ne -2$ וגם $a\ne 1$
\item $a\ne 0$ וגם $a\ne 1$ וגם $a\ne -2$.
\item $\boxed{a\ne 1}$
\item תמיד
\end{multicols}

\end{enumerate}

%%%
%%%%%%
%%%

\noindent\rule{\textwidth-30pt}{0.4pt}


פתרון:
נרצה לדעתי מתי הוקטור המבוקש שייך למרחב השורות, כלומר מתי ניתן להגיע אליו כצירוף לינארי של שורות המטריצה, זה שקול לכך שלממ״ל הבאה יש פתרון
\begin{align*}
   \begin{pmatrix}
    a&1&1&\bigm| &4 \\
    1 & a & 1& \bigm|&-5 \\
    1 & 1 & a& \bigm|& 1 \\
\end{pmatrix} 
&\begin{bmatrix}
    R_2-R_3\rightarrow R_2\\
    R_1-aR_3\rightarrow R_1\\
    R_1\leftrightarrow R_3\\
\end{bmatrix}
\begin{pmatrix}
1 & 1 & a& \bigm|& 1 \\
 0 & a-1 & 1-a& \bigm|&-6 \\
    0&1-a&1-a^2&\bigm| &4-a \\
\end{pmatrix} 
\\&\begin{bmatrix}
    R_3+R_2\rightarrow R_3\\
\end{bmatrix}
\begin{pmatrix}
1 & 1 & a& \bigm|& 1 \\
    0&a-1&1-a&\bigm| &-6 \\
    0 & 0 & 2-a-a^2& \bigm|&-2-a \\
\end{pmatrix} 
\end{align*}
בשביל שהשורה השנייה לא תהיה שורת סתירה חייב להתקיים
$a\neq 1$,
ובנוסף נבדוק מתי
$2-a-a^2=0$:
זה קורה כאשר
$a=1$
או
$a=-2$.
עבור
$a=-2$
נקבל שורת אפסים ולא שורת סתירה בשורה השלישית, ולכן אין כאן בעיה.

לסיכום, הערכים עבורם יהיה פתרון לממ״ל הם 
$a\neq1$.

\noindent\rule{\textwidth-30pt}{0.4pt}


%%%
%%%%%%
%%%
\item\points{5}
תהא
\[A=\begin{pmatrix} 
1 & 2 & 0 & 1 & 3 \\
0 & 1 & 1 & 2 & 1 \\
1 & 3 & 1 & 3 & 4 \\
2 & 5 & 1 & 4 & 7 \\
0 & 2 & 2 & 4 & 2
\end{pmatrix}\]
ותהא
$T_A:\cc^5\to \cc^5$
ההעתקה הלינארית המוגדרת 
על ידי
\ 
$T_A(v) = A\cdot v$.
מה מהבאים
\textbf{
\underline{
לא
}}
נכון?
\begin{enumerate}
\begin{multicols}{2}
    \item $\dim(\row(A))=2$
        \item $\boxed{\dim(N(A))=2}$
    \item $\dim(\im(T_A)) = 2$
    \item $\dim(\row(A^t))=2$
\end{multicols}
\end{enumerate}

\noindent\rule{\textwidth-30pt}{0.4pt}


פתרון:
ראשית, ניתן לשים לב כי
$$\dim(Row(A))=rank(A)$$
וגם
$$\dim(Im(T_A))=\dim(Col(A))=rank(A)$$
וגם
$$\dim(Row(A^t))=\dim(Col(A))=rank(A)$$
כלומר המספרים בסעיפים
א,ג,ד יצאו אותו הדבר, לכן סעיפים אלה חייבים להיות נכונים (אחרת כולם היו לא נכונים), מה שאומר שהתשובה הנכונה היא סעיף ב.

בדרך ישירה, נדרג את
$A$
ונקבל כי
$rank(A)=3$,
לכן
$\dim(N(A))=5-3=2$.

\noindent\rule{\textwidth-30pt}{0.4pt}

\item\points{5}
תהא
$A\in M_{m\times n} (F)$ מטריצה שאינה מטריצת האפס.

אם לממ"ל
$A x=0$
קיים יותר מפתרון אחד, 
אז בהכרח:
\begin{enumerate}
  \item $m>n$.
	\item  לכל $0 \ne b\in F^m$ לממ"ל  $Ax=b$ קיים פתרון.
	\item
   \boxed{\text{קיים
   $0\ne b\in F^m$ כך שלממ"ל  $Ax=b$ יש יותר מפתרון אחד.}}
	\item  אף סעיף אחר לא בהכרח נכון.
\end{enumerate}
\noindent\rule{\textwidth-30pt}{0.4pt}

פתרון: את סעיף (א) נפסול עם הדוגמה
$A=\begin{pmatrix}
     1&0\\
    1&0\\
\end{pmatrix}$
כי
$m=2=n$
ויש ל-$Ax=0$ אינסוף פתרונות.

את סעיף (ב) נוכל לפסול עם אותה
$A$,כי עבור
$b=\begin{pmatrix}
    1\\0\\
\end{pmatrix}$
לממ״ל 
$Ax=b$
אין פתרון.

מהנתון נסיק לאחר דירוג של
$A$
נקבל כי יש לפחות משתנה חפשי אחד, ולכן לכל
$b$
שנקח, לממ״ל
$Ax=b$
לא יהיה פתרון יחיד. מצד שני, לכל
$b$
שנקח עבורה יש פתרון לממ״ל
$Ax=b$
(למשל, כל
$b$
במרחב העמודות של
$A$)
יהיו אוטומטית אינסוף פתרונות, כמבוקש.
לסיום, קיים
$b$
כזה שהוא שונה מאפס, כי המטריצה
$A$
שונה ממטריצת האפס.

%%%
%%%%%%
%%%


\noindent\rule{\textwidth-30pt}{0.4pt}


\item\points{5}
תהי $T:\mathbb{R}^3\to\mathbb{R}^3$ העתקה לינארית חח"ע המקיימת
\[
T\!\begin{pmatrix}1\\0\\1\end{pmatrix}=\begin{pmatrix}-1\\1\\1\end{pmatrix},
\qquad
T\!\begin{pmatrix}1\\2\\3\end{pmatrix}=\begin{pmatrix}-1\\3\\1\end{pmatrix},
\qquad
T\!\begin{pmatrix}-2\\2\\0\end{pmatrix}=\begin{pmatrix}a\\a-2\\ a-4 \end{pmatrix}.
\]
בחרו את הטענה הנכונה:

\begin{enumerate}
\begin{multicols}{2}
    \item $a=1$
        \item $\boxed{a=2}$
    \item $a=3$
    \item לא ניתן לדעת מהנתונים מה ערכו של $a$.
\end{multicols}
\end{enumerate}

פתרון:
נסמן את שלושת הוקטורים עליהם נתון איך פועלת ההעתקה על ידי
$v_1,v_2,v_3$,
כאשר
$v_1$
הוא השמאלי.
נשים לב כי הם תלויים ומתקיים
$v_2-3v_1=v_3$
אז
נשתמש בלינאריות ההעתקה ונקבל
$$\begin{pmatrix}
    2\\0\\-2\\
\end{pmatrix}
=T(v_2)-3T(v_1)=T(v_2-3v_1)=T(v_3)
=\begin{pmatrix}a\\a-2\\ a-4 \end{pmatrix}$$
ולכן
$a=2$.
\end{enumerate}
%%%
%%%%%%
%%%

 \pagebreak

\section*{חלק ב' (60 נק') - יש לבחור 3 מתוך 4 ולסמן בעמוד הראשון.}
יש לענות במקומות המתאימים בעמודים הבאים.
\begin{enumerate}
\newcommand{\firstQ}{
\item \points{20}
תהי
\[A = \begin{pmatrix}
    2 & 2 & 0 & 0\\
    2 & 2 & 0 & 0\\
    0 & 0 & 2 & 2 \\
    0 & 0 & 2 & 2
\end{pmatrix} \]
\begin{enumerate}
    \item\points{5} 
    מצאו את הפולינום האופייני של $A$  ואת כל הערכים העצמיים של $A$.
    \item\points{10}
    קבעו
    האם $A$ לכסינה מעל $\rr$,
    ואם כן, מצאו $D$ אלכסונית
    ו-
    $P$
    הפיכה עבורן
    $A=PDP^{-1}$.
    \item\points{5} 
     חשבו את הפולינום האופייני של המטריצה
     $A^4-4A^3+A-I_4$.
    \end{enumerate}

}
\firstQ
\noindent\rule{\textwidth-30pt}{0.4pt}

\newcommand{\secondQ}{
\item \points{20}
תהי
\[
A=
\begin{pmatrix}
1 & 1 & -1 & 0 \\
0 & 1 & 0  & -1 \\
1 & 1 & -1 & 0 \\
0 & 1 & 0 & -1
\end{pmatrix}
\]
\begin{enumerate}
\item\points{8} מצאו בסיס של $\col(A)$ ובסיס של $N(A)$.
\item\points{8}
מצאו בסיס אורתונורמלי של
$\col(A)+N(A)$.
\item\points{4} מצאו בסיס
אורתונורמלי של
$N(A^2)$.

\end{enumerate}
}
\secondQ

\noindent\rule{\textwidth-30pt}{0.4pt}
\newcommand{\thirdQ}{
\item\points{20}
יהיו 
$B,C$ שני בסיסים סדורים של 
 $\rr^3$ 
  ותהי
  $T:\rr^3\to\rr^3$
     העתקה לינארית כך שמתקיים
    \[
    E=\Set{
    \begin{pmatrix}
    1 \\ 0 \\ 0
    \end{pmatrix},
    \begin{pmatrix}
    0 \\ 1 \\ 0
    \end{pmatrix},
    \begin{pmatrix}
    0 \\ 0 \\ 1
    \end{pmatrix}
    }
    \text{\quad עבור \quad    }
    [T]^{B}_{C} = \begin{pmatrix}
        1 & 2 & 3 \\
        2 & 4 & 6 \\
        1 & 1 & 1 \\
    \end{pmatrix}, [I]^{E}_C = \begin{pmatrix}
    1 & 3 & 5 \\
    0 & 1 & 4 \\
    0 & 0 & 1
    \end{pmatrix}
    \]
    
    \begin{enumerate}
    \item\points{10} מצאו את $[T]^{B}_E$.
    \item\points{5} מצאו בסיס ל-
    $\im(T)$.
    \item\points{2}
    הראו שקיים 
    $w\in \rr^3$
    כך שמתקיים
    $\Ker(T) = \Span(w)$.
    \item\points{3} עבור $w$ כלשהו המקיים
    $\Ker(T) = \Span(w)$,
    מצאו את
    וקטור הקואורדינטות
    $[w]_B$.

    (אין צורך למצוא את 
    $w$
    עצמו)
    \end{enumerate}
}
\thirdQ

\noindent\rule{\textwidth-30pt}{0.4pt}
\newcommand{\forthQ}{
    \item \points{20}
    הסעיפים הבאים לא קשורים אחד לשני:
    \begin{enumerate}
        \item\points{10} 
        יהיו
        $v_1,v_2,v_3\in \rr^5$ ותהי $A\in M_{5\times 5}(\rr)$.

        הראו כי אם הוקטורים
        $v_1,v_2,v_3$ בלתי תלויים לינארית (בת"ל) ואם המטריצה $A$ היא הפיכה,
        
        אז
        הוקטורים
        $Av_1,Av_2,Av_3$
        הם בלתי תלויים לינארית.

        \ 
        
        \item\points{10} 
        נניח ש-
        $A+B = I_n$ וגם
        $\lambda=1$
        הוא ערך עצמי של המטריצה
        $B$.
        הוכיחו כי $A$ לא הפיכה.
    \end{enumerate}
}\forthQ
\pagebreak
\setcounter{enumi}{0}
    \firstQ
\noindent\rule{\textwidth-30pt}{0.4pt}
\underline{פיתרון:}
\begin{enumerate}
    \item 
    נחשב את הפולינום האופייני של המטריצה
    \begin{align*}
        &p_A(x)=\det(xI_4-A)=
        \det\begin{pmatrix}
            x-2&-2&0&0\\
            -2&x-2&0&0\\
            0&0&x-2&-2\\
            0&0&-2&x-2\\
        \end{pmatrix}
        \underset{\begin{bmatrix}
        R_1-R_2\rightarrow R_1\\
        R_3-R_4\rightarrow R_3
        \end{bmatrix}}{=}
        \\&
        \det\begin{pmatrix}
            x&-x&0&0\\
            -2&x-2&0&0\\
            0&0&x&-x\\
            0&0&-2&x-2\\
        \end{pmatrix}
        \underset{\begin{bmatrix}
        R_1-R_2\rightarrow R_1\\
        R_3-R_4\rightarrow R_3
        \end{bmatrix}}{=}x^2   
          \det\begin{pmatrix}
            1&-1&0&0\\
            -2&x-2&0&0\\
            0&0&1&-1\\
            0&0&-2&x-2\\
        \end{pmatrix}\\&
        \underset{\begin{bmatrix}
        R_2+2R_1\rightarrow R_2\\
        R_4+2R_3\rightarrow R_4\\
        \end{bmatrix}}{=}
           \det\begin{pmatrix}
            1&-1&0&0\\
            0&x-4&0&0\\
            0&0&1&-1\\
            0&0&0&x-4\\
        \end{pmatrix}=x^2(x-4)^2
    \end{align*}
    כאשר השלב האחרון נובע כי המטריצה היא משולשית עליונה ולכן הדטרמיננטה שווה למכפלת איברי האלכסון. מכך ניתן לראות כי הערכים העצמיים של
    $A$
    הם
    $\lambda=0,4$.
    
    \item 
    בשביל לבדוק האם המטריצה לכסינה, נמצא את הריבויים האלגברי והגיאומטרי של כל ערך עצמי. עבור הערך העצמי
    $\lambda=0$,
    קיבלנו כי הריבוי האלגברי הוא
    $2$,
    ואת הריבוי הגיאומטרי נמצא לאחר 
    מציאת המרחב העצמי המתאים
    $N(0\cdot I_4-A)=N(-A)=N(A)$,
    על ידי דירוג המטריצה
    \begin{align*}
     A
     \begin{bmatrix}
         R_2-R_1\rightarrow R_2\\
         R_4-R_3\rightarrow R_4\\
         1/2R_1\rightarrow R_1\\
         1/2R_3\rightarrow R_3\\
         R_2\leftrightarrow R_2\\
     \end{bmatrix}
     \begin{pmatrix}
         1&1&0&0\\
         0&0&1&1\\
         0&0&0&0\\
         0&0&0&0\\
     \end{pmatrix}
    \end{align*}
    לכן הפתרון הכללי של הממ״ל נתון על ידי
    $$ (-s , s,-t , t)=s(-1,1,0,0)+t(0,0,-1,1)$$
    לכן 
    $\dim(N(A))=2$
    וזהו הריבוי הגיאומטרי של הערך העצמי
    $\lambda=0$.

עבור הערך העצמי
    $\lambda=4$,
    קיבלנו כי הריבוי האלגברי הוא
    $2$,
    ואת הריבוי הגיאומטרי נמצא לאחר 
    מציאת המרחב העצמי המתאים
    $N(4\cdot I_4-A)$,
    על ידי דירוג המטריצה
    \begin{align*}
     4I_4-A=\begin{pmatrix}
            2&-2&0&0\\
            -2&2&0&0\\
            0&0&2&-2\\
            0&0&-2&2\\
        \end{pmatrix}
     \begin{bmatrix}
         R_2+R_1\rightarrow R_2\\
         R_4+R_3\rightarrow R_4\\
         1/2R_1\rightarrow R_1\\
         1/2R_3\rightarrow R_3\\
         R_2\leftrightarrow R_2\\
     \end{bmatrix}
     \begin{pmatrix}
         1&-1&0&0\\
         0&0&1&-1\\
         0&0&0&0\\
         0&0&0&0\\
     \end{pmatrix}
    \end{align*}
    לכן הפתרון הכללי של הממ״ל נתון על ידי
    $$ (s , s,t,t) =s(1,1,0,0)+t(0,0,1,1)$$
    לכן 
    $\dim(N(4I_4-A))=2$
    וזהו הריבוי הגיאומטרי של הערך העצמי
    $\lambda=4$. 

    לסיכום, קיבלנו כי המטריצה
    $A$
    לכסינה מעל
    $\rr$
    ומתקיים
    $A=PDP^{-1}$
    כאשר
    $$.D=
    \begin{pmatrix}
    4&0&0&0\\
    0&4&0&0\\
    0&0&0&0\\
    0&0&0&0\\
    \end{pmatrix},\quad
    P=\begin{pmatrix}
        1&0&1&0\\
        1&0&-1&0\\
        0&1&0&1\\
        0&1&0&-1\\
    \end{pmatrix}$$
    \item 
    נסמן את המטריצה
    $B=A^4-4A^3+A-I_4$
    וניעזר בצורה המיוחדת של
    $A$
    שנבעה מההליכסון
    \begin{align*}
       B&=(PDP^{-1})^4-4(PDP^{-1})^3+PDP^{-1}-I_4
       \\&=P(D^4-4D^3+D-I_4)P^{-1}
       =P \begin{pmatrix}
           3&0&0&0\\
           0&3&0&0\\
           0&0&-1&0\\
           0&0&0&-1\\
       \end{pmatrix}P^{-1}
    \end{align*}
    ולכן
    $B$
    דומה למטריצה האלכסונה האחרונה שרשמנו, אז היא לכסינה והפולינום האופייני שלה שווה
    $$.p_B(x)=(x+1)^2(x-3)^2$$
\end{enumerate}    
    \pagebreak \ 
    


    
    \secondQ
    \noindent\rule{\textwidth-30pt}{0.4pt}
    \underline{פיתרון:}
\begin{enumerate}
    \item 
ראשית, נדרג את המטריצה ונקבל
\begin{align*}
    A
    \begin{bmatrix}
    R_3-R_1\rightarrow R_3\\
    R_4-R_2\rightarrow R_4\\
    R_1-R_2\rightarrow R_1\\
    \end{bmatrix}
    \begin{pmatrix}
        1&0&-1&1\\
        0&1&0&-1\\
        0&0&0&0\\
        0&0&0&0\\
    \end{pmatrix}
\end{align*}
בגלל שקיבלנו איבר מוביל רק בעמודות הראשונה והשנייה, נסיק כי בסיס למרחב העמודות יהיה הקבוצה שמורכבת מהוקטור הראשון והוקטור השני במקור, כלומר
$$.B_{Col(A)}=
\left\{
v_1=(1,0,1,0),
v_2=(1,1,1,1)\right\}$$
למציאת בסיס עבור
$N(A)$,
נפתור את הממ״ל: 
מהדירוג קבלנו כי
$x_3=t$
וגם
$x_4=s$
משתנים חפשיים, ובנוסף
$x_1=t-s, x_2=s$
לכן מרחב הפתרונות יהיה
\begin{align*}
    N(A)=
    \left\{(
        t-s,s,t,s)
  : t,s\in\rr
    \right\}=\Span
    \left(
v_1=(1,0,1,0),
v_3=(-1,1,0,1)
    \right)
\end{align*}
כאשר שני הוקטורים האחרונים הם בת״ל ולכן בסיס של
$N(A)$.
    \item 
מהסעיף הקודם
$\{v_1,v_2\}$
בסיס של
$Col(A)$,
ו-$\{v_1,v_3\}$
בסיס של
$N(A)$,
לכן נקבל כי
$$Col(A)+N(A)=Span(v_1,v_2)+Span(v_1,v_3)=Span(v_1,v_2,v_3)$$
אז כל שנותר הוא למצוא בסיס לפרישה האחרונה. נעשה זאת למשל על ידי כך שנשים את שלושת הוקטורים בשורות מטריצה ונדרג אותה
\begin{align*}
\begin{pmatrix}
    1&0&1&0\\
    1&1&1&1\\
    -1&1&0&1\\
    \end{pmatrix}
&\begin{bmatrix}
    R_2-R_1\rightarrow R_2\\
    R_3+R_1\rightarrow R_3\\
\end{bmatrix}
\begin{pmatrix}
    1&0&1&0\\
    0&1&0&1\\
    0&1&1&1\\
    \end{pmatrix}
\\&\begin{bmatrix}
    R_3-R_2\rightarrow R_3\\
\end{bmatrix}
\begin{pmatrix}
    1&0&1&0\\
    0&1&0&1\\
    0&0&1&0\\
    \end{pmatrix}
    \begin{bmatrix}
    R_1-R_3\rightarrow R_1\\
\end{bmatrix}
\begin{pmatrix}
    1&0&0&0\\
    0&1&0&1\\
    0&0&1&0\\
    \end{pmatrix}
\end{align*}
ולכן בסיס עבור
$Col(a)+N(A)$
יהיה מורכב משלושת השורות שקיבלנו לאחר הדירוג (כי אין שורות אפסים), כלומר
$$B=\left\{
(1,0,0,0),
(0,1,0,1),
(0,0,1,0)
\right\}$$
במקרה נשים לב כי זהו בסיס אורתוגונלי, ולכן כל שנשאר זה לנרמל את הקוטורים שבו. הראשון והאחרון הם כבר מנורמלים, אז רק נחלק את הוקטור השני בנורמה (אורך) שלו שיוצא
$\sqrt{2}$
ונקבל את הבסיס האורתונורמלי
$$.B'=\left\{
(1,0,0,0),
\frac{1}{\sqrt{2}}(0,1,0,1),
(0,0,1,0)
\right\}$$

$\bullet$
\underline{הערה ראשונה:}
ניתן היה לשים את הוקטורים בעמודות כדי לוודא שהם בת״ל ואז היינו מקבלים את שלושת הוקטורים
$v_1,v_2,v_3$
כבסיס, ובמצעים עליו את תהליך גראם-שמידט, כדי לקבל בסיס אורתונורמלי, שהיה יוצא למשל
$$\left\{
\frac{1}{\sqrt{2}}(1,0,1,0),\, \frac{1}{\sqrt{2}}(0,1,0,1),\,
\frac{1}{\sqrt{2}}(-1,0,1,0)
\right\}$$

$\bullet$
\underline{הערה שנייה:}
שימו לב כי אם מבצעים גראם-שמידט על הוקטורים בסדר אחר, יכול להתקבל בסיס אורתונורמלי שנראה אחרת, וזה בסדר גמור. אין כאן תשובה יחידה- יש אינסוף בסיסים אורתונורמלים המהווים תשובה לשאלה.
    \item 
    בשביל למצוא בסיס אורתונורמלי למרחב הפתרונות של
    $A^2$,
    ראשית נחשב את המטריצה הזו:
    $$A^2=
    \begin{pmatrix}
        0&1&0&-1\\
        0&0&0&0\\
        0&1&0&-1\\
        0&0&0&0\\
    \end{pmatrix}$$
    אחרי דירוג נישאר רק עם שורה אחת שאינה שורת אפסים המתארת את המשוואה
    $x_2-x_4=0$,
    אז
    \begin{align*}
        N(A^2)=\left\{
(a,d,c,d): a,c,d\in\rr
        \right\}=\Span
        \left\{
(1,0,0,0),\, (0,1,0,1),\, (0,0,1,0)
\right\}
    \end{align*}
    נשים לב כי זו אותה קבוצת וקטורים שקיבלנו בסעיף הקודם, כלומר
    $$N(A^2)=Col(a)+N(A)$$
    ולכן בסיס אורתונורמלי יהיה זהה לתשובה שמצאנו בסעיף הקודם.
\end{enumerate}
\ \pagebreak

    \thirdQ
    \noindent\rule{\textwidth-30pt}{0.4pt}
    \underline{פיתרון:}
\begin{enumerate}
    \item 
    ניזכר כי בכיתה למדנו נסוחה שאומרת כי 
    $$[T]_E^B=[I]^C_E\cdot [T]_C^B\cdot [I]_B^B=[I]^C_E\cdot [T]^B_C$$
    ובנוסף
    $[I]_E^C=([I]_C^E)^{-1}$
    ולכן נחשב את ההופכית של המטריצה הנתונה בעזרת האלגוריתם למציאת הופכית
    \begin{align*}
\begin{pmatrix}
    1&3&5&\bigm|&1&0&0\\
    0&1&4&\bigm|&0&1&0\\
    0&0&1&\bigm|&0&0&1\\
\end{pmatrix}
\begin{bmatrix}
    R_2-4R_3\rightarrow R_2\\
    R_1-5R_3\rightarrow R_1\\
    R_1-3R_2\rightarrow R_2\\
\end{bmatrix}
\begin{pmatrix}
    1&0&0&\bigm|&1&-3&7\\
    0&1&0&\bigm|&0&1&-4\\
    0&0&1&\bigm|&0&0&1\\
\end{pmatrix}
    \end{align*}
    לכן
    $$.[T]_E^B=
    \begin{pmatrix}
      1&-3&7\\
      0&1&-4\\
      0&0&1\\
    \end{pmatrix}\cdot
     \begin{pmatrix}
        1 & 2 & 3 \\
        2 & 4 & 6 \\
        1 & 1 & 1 \\
    \end{pmatrix}
    =\begin{pmatrix}
    2&-3&-8\\
    -2&0&2\\
    1&1&1\\
    \end{pmatrix}
    $$
    \item 
    נסמן
    $B=\{v_1,v_2,v_3\}$,
ומההגדרה של מטריצה מייצגת
(אותה מצאנו בסעיף הקודם)
נקבל
$$T(v_1)=\begin{pmatrix}
    2\\-2\\1\\
\end{pmatrix},\quad T(v_2)=\begin{pmatrix}
    -3\\0\\1\\
\end{pmatrix},\quad T(v_3)=\begin{pmatrix}
    -8\\2\\1\\
\end{pmatrix}$$
ובגלל ששלושת הוקטורים האלו מהווים בסיס של
$\rr^3$
נקבל כי
התמונה של
$T$
נפרשת על ידי שלושת הוקטורים הללו. כדי להבין אם הם בסיס או לא, נשים אותם בעמודות, כלומר נסתכל על המטריצה
$[T]_E^B$
ונדרג
\begin{align*}\begin{pmatrix}
    2&-3&-8\\
    -2&0&2\\
    1&1&1\\
\end{pmatrix}
\begin{bmatrix}
R_2+R_1\rightarrow R_2\\
    R_1-2R_3\rightarrow R_1\\
    R_1\leftrightarrow R_3\\
\end{bmatrix}
\begin{pmatrix}
1&1&1\\
    0&-3&-6\\
        0&-5&-10\\
\end{pmatrix}
\begin{bmatrix}
-1/3R_2\rightarrow R_2\\
    R_1-R_2\rightarrow R_1\\
    R_3+5R_2\leftrightarrow R_3\\
\end{bmatrix}
\begin{pmatrix}
1&0&-1\\
    0&1&2\\
        0&0&0\\
\end{pmatrix}
\end{align*}
קיבלנו איברים מובילים בעמודות הראשונה והשנייה, לכן מלכתחילה שני הוקטורים הראשונים יהיו בסיס של התמונה, כלומר הקבוצה
$$\left\{
\begin{pmatrix}
    2\\-2\\1\\
\end{pmatrix},\begin{pmatrix}
    -3\\0\\1\\
\end{pmatrix}
\right\}$$
בסיס של
$\Img(T)$.
    \item 
מסעיף קודם קיבלנו כי
$\dim(\Img(T))=2$,
ולכן ממשפט המימדים של העתקות לינאריות נסיק
$$\dim(\ker(T))=3-\dim(\Img(T))=1$$
מה שאומר שקיים
לגרעין בסיס עם וקטור אחד בדיוק
$w\in\rr^3$,
ואז
$\ker(T)=\Span(w)$.
    \item 
   אנו לא יודעים מיהו אותו
   $w$,
   אבל אם
   $\ker(T)=\Span(w)$
   אז
   $w\in\ker(T)$,
   כלומר
   $T(w)=0$.
   נקשר את זה למטריצה המייצגת של ההעתקה בעזרת הנוסחה הבאה
   $$[T(w)]_E=[T]_E^B\cdot [w]_B$$
   אבל אגף שמאל שווה לוקטור האפס, לכן
   $[w]_B$
   יהיה פתרון של הממ״ל ההומוגנית של המטריצה הנתונה מסעיף הראשון. דירגנו אותה בסעיף השני, ולכן כל הפתרונות של הממ״ל ההומוגנית שלה יהיו מהצורה
   $$ 
       ,(t,-2t,t)
$$
   אז תשובה אפשרית תהיה
   $(1,-2,1)$
   או
   $(3,-6,3)$,
   או כל וקטור מהצורה 
   $(t,-2t,t)$
   עבור
   $t\neq0$
   יעבוד ותייצר פיתרון נכון.
\end{enumerate}

  \ \pagebreak

    \forthQ
    \noindent\rule{\textwidth-30pt}{0.4pt}
    \underline{פיתרון:}
\begin{enumerate}
    \item 
    נרצה להראות כי
    $Av_1,Av_2,Av_3$
    בת״ל, לכן נניח כי קיימים
    $\beta_1,\beta_2,\beta_3\in\rr$
    עבורם מתקיים
    $$\beta_1\cdot Av_1+\beta_2\cdot Av_2+\beta_3\cdot Av_3=0$$
    ונסיים את ההוכחה בכך שנראה כי
    $\beta_1=\beta_2=\beta_3=0$
    בהכרח.
    נפשט את מה שכתוב מעלה ונקבל
    $$A\cdot(\beta_1v_1+\beta_2v_2+\beta_3v_3)=0$$
    אך
    $A$
    הפיכה ולכן לממ״ל ההומוגנית שלה יש רק את הפתרון הטריוויאלי, אז בהכרח
    $\beta_1v_1+\beta_2v_2+\beta_3v_3=0$.
    לסיום, מהנתון כי
    $v_1,v_2,v_3$
    בת״ל נסיק כי בהכרח
    $\beta_1=\beta_2=\beta_3=0$,
    כמבוקש.
\item[] 
    \item 
    מהנתון כי
    $\lambda=1$
    הוא ערך עצמי של
    $B$
    נסיק כי קיים
    $v\neq 0$
    עבורו מתקיים
    $Bv=1\cdot v=v$.
    נשתמש בנתון הראשון, נכפיל את שני האגפים בוקטור
    $v$
    מימין ונקבל
    $$Av+v=Av+Bv=(A+B)v=I_n\cdot v=v$$
    לכן לאחר העברת אגף נקבל
    $Av=0$.
    נזכיר כי
    $v\neq0$,
    ולכן לממ״ל ההומוגנית 
    $AX=0$
    יש פתרון לא טריוויאלי, כלומר
    $A$
    לא הפיכה.
\end{enumerate}

    \pagebreak \ \pagebreak
\end{enumerate}

    \pagebreak 
    
    עמוד נוסף לכתיבת המשך פיתרונות
    
    \pagebreak
    עמוד נוסף לכתיבת המשך פיתרונות

 \pagebreak

 % --- דף הנוסחאות ---

\section*{דף הנוסחאות}

\begin{enumerate}
\item $M_{m\times n}(F)$ זו קבוצת המטריצות בגודל $n\times m$
($m$ שורות, $n$ עמודות)
עם רכיבים מ-$F$.

\noindent\rule{\textwidth-30pt}{0.4pt}

\item לכל מטריצה $A$, נסמן על ידי $(A)_{i,j}$ את האיבר של $A$ הנמצא בשורה ה-$i$ ובעמודה ה-$j$.

\noindent\rule{\textwidth-30pt}{0.4pt}

\item פיתוח דטרמיננטה של $A\in M_{n\times n}(F)$ לפי השורה ה-$i$ או לפי העמודה ה-$j$:
\[
\det(A)=\sum_{i=1}^{n}(-1)^{i+j}a_{ij}\det\!\big(M_{ij}(A)\big)
=\sum_{j=1}^{n}(-1)^{i+j}a_{ij}\det\!\big(M_{ij}(A)\big)
\]
כאן $M_{ij}(A)$ זה המינור ה-$i,j$ של המטריצה $A$.

\noindent\rule{\textwidth-30pt}{0.4pt}

\item אם $T:V\to W$ העתקה לינארית, אם $B=\{v_1,\ldots,v_n\}$ בסיס סדור של $V$ ואם $C$ בסיס סדור של $W$,
אז המטריצה המייצגת של $T$ ביחס לבסיסים $B$ ו-$C$ מוגדרת על ידי
\[
[T]^B_C=
\begin{pmatrix}
\vline & & \vline\\
[T(v_1)]_C & \cdots & [T(v_n)]_C\\
\vline & & \vline
\end{pmatrix}.
\]

\noindent\rule{\textwidth-30pt}{0.4pt}

\item לכל $u=(a_1,\ldots,a_n),\; w=(b_1,\ldots,b_n)\in\mathbb{R}^n$

נסמן את המכפלה הפנימית הסטנדרטית ב-$\mathbb{R}^n$ (המכפלה הסקלרית):
\[
 u\cdot w=a_1\cdot b_1+\cdots+a_n\cdot b_n.
\]

\noindent\rule{\textwidth-30pt}{0.4pt}

% --- דף הנוסחאות ---


\item תהליך גראם-שמיט (גרסה א):
אם $W$ תת-מרחב של $\rr^n$, ואם $\{v_1,\ldots,v_m\}$ בסיס של $W$

אז $\{w_1,\ldots,w_m\}$ בסיס אורתוגונלי של $W$ כאשר
\[
w_1=v_1,\qquad
w_k=v_k-\sum_{i=1}^{k-1}\frac{ v_k\cdot w_i}{ w_i\cdot w_i}\,w_i
\quad (k=2,\ldots,m).
\]
\item תהליך גראם-שמיט (גרסה ב):
בהנתן סדרה בת"ל
$\{v_1,\ldots,v_k\}\subseteq \rr^n$, נבנה סדרה אורתונורמלית באופן הבא:
\[
w_1=\frac{v_1}{\|v_1\|},\qquad
w_i'=v_i-\sum_{l=1}^{i-1} (v_i\cdot w_l)\,w_l,\qquad w_i = \frac{w_i'}{\|w_i\|}
\]

\noindent\rule{\textwidth-30pt}{0.4pt}

\item השורשים של המשוואה הריבועית
$
0=c+bx+ax^2,\qquad a\ne 0,
$
נתונים על ידי
\[
x_{1,2}=\frac{-b\pm\sqrt{b^2-4ac}}{2a}.
\]

\noindent\rule{\textwidth-30pt}{0.4pt}

\item מטריצה $A\in M_{n\times n}(\ff)$ דומה למטריצה $B\in M_{n\times n}(\ff)$ 

אם
קיימת 
מטריצה הפיכה
$P\in M_{n\times n}(\ff)$ כך ש-
$A=PBP^{-1}$.
\end{enumerate}

\end{document}